\subsection{Überblick zu KI in ERP-Systemen}
\ac{KI} hat sich in den letzten Jahren zu einer disruptiven Technologie entwickelt, die die Art und Weise verändert, wie verschiedenste Branchen und Geschäftsprozesse funktionieren. Im Kontext von \ac{ERP}-Systemen eröffnet der Einsatz von \ac{KI} durch die Ermöglichung der Analyse riesiger Datenmengen, die Erkennung von Trends und der Erleichterung von Entscheidungsprozessen neue Möglichkeiten zur Optimierung und Automatisierung betrieblicher Abläufe. \cite{fathima_impact_2024} Dabei lassen sich unterschiedliche Einsatzformen von \ac{KI} in ERP-Systemen identifizieren, darunter insbesondere Predictive Analytics, generative \ac{KI} sowie \ac{KI}-Agenten.

\subsubsection{Predictive Analytics}
Im Bereich der Predictive Analytics wird \ac{KI} in \ac{ERP}-Systemen eingesetzt, um zukünftige Entwicklungen auf Basis historischer und aktueller Unternehmensdaten hervorzusagen. Durch den Einsatz von Machine-Learning-Algorithmen können hierfür größere Datenmengen analysiert, komplexe Muster erkannt und präzisere Vorhersagen getroffen werden, als mit traditionellen Forecasting-Methoden. Die Integration \ac{KI}-basierter Predictive Analytics in ERP-Systeme ermöglicht es Unternehmen, Nachfrageentwicklungen frühzeitig zu erkennen, Bestände zu optimieren und begründete Entscheidungen zu treffen. Darüber hinaus unterstützen Echtzeitdaten und kontinuierlich aktualisierte Prognosemodelle eine höhere Reaktionsfähigkeit auf dynamische Marktveränderungen, was insbesondere die Effizienz von Lieferketten und die Ressourcennutzung verbessert. Insgesamt trägt der Einsatz von KI-gestützter Predictive Analytics dazu bei, Prognosegenauigkeit, operative Effizienz und Entscheidungsqualität in ERP-Systemen signifikant zu steigern. \cite{fathima_impact_2024}

\subsubsection{Generative KI}
Generative \ac{KI} eröffnet eine Vielzahl von Anwendungsfällen in \ac{ERP}-Systemen, die sich über nahezu alle Module erstrecken. Ein grundlegender Einsatzbereich liegt in der Unterstützung bei der Erstellung textbasierter Inhalte. Dazu zählen unter anderem die automatische Generierung von E-Mails, Marketing- und Medieninhalten sowie die Erstellung von Produktbeschreibungen. \cite[\pagef 73346]{sarferaz_implementing_2025}

Darüber hinaus kann generative \ac{KI} zur Beantwortung von Nutzeranfragen sowie in Form von Konversationsagenten innerhalb von \ac{ERP}-Systemen eingesetzt werden. Diese unterstützen Nutzer bei der Navigation durch komplexe Prozesse, stellen kontextbezogene Informationen bereit und können Handlungsempfehlungen auf Basis vorhandener Systemdaten geben. Dadurch wird der Benutzersupport verbessert, die Einarbeitungszeit neuer Nutzer reduziert und die Effizienz bei der Systemnutzung gesteigert. \cite[\pagef 73346]{sarferaz_implementing_2025}

Weitere Einsatzbereiche generativer \ac{KI} betreffen die Verarbeitung und Aufbereitung von Informationen in \ac{ERP}-Systemen. Dazu zählen sowohl die automatische Zusammenfassung umfangreicher textbasierter Inhalte als auch die Extraktion strukturierter Informationen aus unstrukturierten Dokumenten, etwa aus Rechnungen, Lieferdokumenten oder Verträgen. Dadurch können relevante Inhalte schneller erfasst, manuelle Dateneingaben reduziert und Prozesse effizienter gestaltet werden. \cite[\pagef 73347]{sarferaz_implementing_2025}

Des Weiteren kann generative \ac{KI} zur Unterstützung komplexer Konfigurations- und Entscheidungsprozesse eingesetzt werden. In diesem Kontext ist sie in der Lage, fehlende oder inkonsistente Systemeinstellungen zu identifizieren und Vorschläge zur Ergänzung oder Korrektur von Konfigurationen zu generieren. Die abschließende Prüfung und Übernahme der vorgeschlagenen Maßnahmen erfolgt dabei weiterhin durch den Anwender, wodurch die Kontrolle über geschäftskritische Prozesse gewährleistet bleibt. \cite[\pagef 73347]{sarferaz_implementing_2025}

Zusammenfassend zeigen die genannten Anwendungsfälle, dass generative \ac{KI} in \ac{ERP}-Systemen insbesondere Chancen in der Reduktion manueller Tätigkeiten bietet. Durch die Automatisierung wissensintensiver Aufgaben, die Unterstützung bei der Informationsaufbereitung sowie die verbesserte Interaktion zwischen Nutzer und System kann die Produktivität signifikant erhöht werden.

\subsubsection{KI-Agenten}
Während generative \ac{KI} vor allem auf die Erstellung und Verarbeitung von Inhalten fokussiert ist und dabei eine assistierende Rolle einnimmt, gehen \ac{KI}-Agenten in \ac{ERP}-Systemen einen Schritt weiter, indem sie eigenständige Handlungen ausführen können. Durch die Integration von \ac{KI}-Agenten in \ac{ERP}-Systeme kann die Automatisierung vorangetrieben werden, da diese Agenten in der Lage sind, Entscheidungen in Echtzeit zu treffen und selbstständig zu lernen. Diese Fähigkeiten ermöglichen dadurch die Optimierung von Arbeitsabläufen sowie eine Verbesserung der Produktivität und der Kundenerfahrungen. \cite[\pagef 1032]{gujar_role_2025}

Ein Anwendungsfall von \ac{KI}-Agenten in \ac{ERP}-Systemen ist, genau wie bei der generativen \ac{KI}, \ac{KI}-gestützte Chatbots. Diese Chatbots können allerdings, anders als rein generative \ac{KI}, nicht nur auf Anfragen mit Antworten reagieren, sondern auch eigenständig Aktionen im \ac{ERP}-System durchführen. Dadurch können Supportanfragen nicht nur beantwortet, sondern identifizierte Probleme direkt bearbeitet werden. \cite[\pagef 1033]{gujar_role_2025} 

Ein weiterer Anwendungsfall in der Kundenbetreuung ist die Nutzung von \ac{KI}-Agenten zur proaktiven Identifikation und Lösung von Kundenproblemen. Durch die Analyse von Kundendaten und Interaktionen können \ac{KI}-Agenten potenzielle Probleme, wie beispielsweise Lieferverzögerungen, frühzeitig erkennen und entweder den Kunden informieren oder automatisch Maßnahmen zur Problemlösung einleiten. \cite[\pagef 1033]{gujar_role_2025}

Neben der Kundenbetreuung können \ac{KI}-Agenten auch für die Automatisierung und Optimierung von Geschäftsprozessen in \ac{ERP}-Systemen eingesetzt werden, indem sie wiederkehrende Workflows eigenständig ausführen und dabei mit Daten aus unterschiedlichen Systemen arbeiten. Insbesondere Prozesse wie die Auftragsabwicklung, das Bestandsmanagement oder das Onboarding von Kunden lassen sich durch \ac{KI}-Agenten automatisieren, sodass Aufgaben effizienter und fehlerärmer durchgeführt werden können. \cite[\pagef 1034]{gujar_role_2025}

Zusammenfassend zeigen die dargestellten Einsatzmöglichkeiten von \ac{KI}-Agenten in \ac{ERP}-Systemen, dass sich erhebliche Effizienzpotenziale ergeben. Durch die Automatisierung wiederkehrender und wissensintensiver Aufgaben können Prozessdurchlaufzeiten verkürzt und manuelle Tätigkeiten reduziert werden. Darüber hinaus ermöglichen \ac{KI}-Agenten eine proaktive Identifikation und Bearbeitung von Problemen, wodurch Serviceprozesse beschleunigt, Mitarbeitende entlastet und die Servicequalität verbessert werden können. \cite[\pagef 1032ff]{gujar_role_2025} Die Fähigkeit zur systemübergreifenden Interaktion erweitert diese Potenziale zusätzlich, da \ac{KI}-Agenten Daten aus verschiedenen Unternehmenssystemen integrieren und Entscheidungen auf einer umfassenderen Datenbasis treffen können. \cite[\pagef 8f]{sarferaz_implementing_2025-1}

%Ein wesentlicher Vorteil von \ac{KI}-Agenten in \ac{ERP}-Systemen liegt also auch in der Fähigkeit, nicht nur im \ac{ERP}-System selbst zu agieren, sondern auch systemübergreifend zu arbeiten. Dadurch können diese Agenten Daten aus verschiedenen Unternehmenssystemen integrieren und so umfassendere Entscheidungen treffen. \cite[\pagef 8f]{sarferaz_implementing_2025-1}